%% 70-sustainability.tex — Sustainability

\section{Sustainability}
\label{sec:sustainability}

\paragraph{Versioning.}
Per-item provenance is recorded at ingest time using the W3C PROV Ontology~\cite{lebo2013provo}: each \texttt{edm:ProvidedCHO} carries \texttt{prov:generatedAtTime}, \texttt{dcterms:hasVersion} (DDB revision id), and links to the source dataset and mapping-version agent derived from the DDB JSON metadata.
Graph-level pipeline provenance --- an Activity record per script run, capturing agent identity, version, timestamps, and an \texttt{inferenceMethod} flag (\texttt{Deterministic} / \texttt{Heuristic} / \texttt{NER}) --- is planned as future work.

\paragraph{Update cadence.}
GEMEA is updated periodically to track DDB data releases.

\paragraph{Persistent URI.}
The dataset landing page at \url{https://gemea.ise.fiz-karlsruhe.de/} serves as the stable reference URI, with downloads at \url{https://gemea.ise.fiz-karlsruhe.de/downloads/} and the SHMARQL browser at \url{https://gemea.ise.fiz-karlsruhe.de/shmarql/}.

\paragraph{Open source.}
The GEMEA index is reproducible from the DDB data dump using the pipeline scripts in the repository.

\paragraph{Institutional support.}
GEMEA is developed at FIZ Karlsruhe -- Leibniz Institute for Information Infrastructure in collaboration with the Karlsruhe Institute of Technology (AIFB).
The resource is hosted on FIZ Karlsruhe infrastructure and maintained indefinitely as part of the institute's research data services.

\paragraph{Dependency management.}
The \textbf{MOCHO} alignment layer depends on \texttt{mocho.owl}, which is under active development.
Alignment rule changes upstream are tracked in the GEMEA repository and trigger a re-run of the affected named graph.
Known scope boundaries in the current alignment (Section~\ref{sec:quality}) are documented as versioned issue records.

\paragraph{Outbound URI persistence.}
%Outbound links encoded as \texttt{edm:isShownAt} and \texttt{edm:object} values point to \texttt{ddb.de/item/\{id\}} URIs.
As DDB does not operate under a formal persistent URI policy, objects withdrawn or re-harvested by contributing institutions will render a proportion of these links unreachable over time.
